\section{Frozen coefficient and two exact selectors}
\label{sec:selectors}

\subsection{Wordwise and empty-word conventions}

Let \(A_m^+\) denote the positive words in \(A_m\).  The support of a word is
the set of letters that occur in it.

\begin{definition}[pure-power cyclic coefficient]
For \(w\in A_m^+\), define
\begin{equation}
 \chi_m(w)=\mathbf 1_{\{|\supp(w)|=1\}}.
 \label{eq:chi-definition}
\end{equation}
The equality in \eqref{eq:chi-definition} is imposed on each individual
positive word before any commuting color variables are introduced.
\end{definition}

Cyclic rotation does not change support, and
\(\supp(w^r)=\supp(w)\).  Thus \(\chi_m\) is an exact cyclic and
all-repetition coefficient.  The empty word needs a separate convention.

\begin{convention}[empty-word firewall]
We use two extensions to \(A_m^*=A_m^+\cup\{\epsword\}\).
\begin{enumerate}
 \item The \emph{character convention} sets
 \(\chi_m(\epsword)=m\), the identity trace on the canonical color fiber.
 \item The \emph{language convention} sets
 \(\chi_m(\epsword)=0\), so the series is the characteristic series of the
 nonempty pure-power language.
\end{enumerate}
The empty word never occurs in a determinant trace logarithm, but it changes
minimal linear realization by one dormant mode.
\end{convention}

This distinction also separates two theorem classes.  A general recognizable
series may have a bilinear readout
\(f(w)=\alpha\mu(w)\beta\); its minimal dimension is controlled by the Hankel
matrix.  The character theorem in \Cref{sec:character} assumes instead that
the coefficient is a matrix trace or supertrace of a stationary
multiplicative representation.  No character classification is asserted for
an arbitrary bilinear readout.

\subsection{An orbitwise support-incidence coefficient}

For a nonempty finite set \(S\), define its reduced permutation space
\begin{equation}
 Q_S=\ker\!\left(\C^S\xrightarrow{\sum}\C\right),
 \qquad \dim Q_S=|S|-1.
 \label{eq:reduced-support}
\end{equation}
Give \(\Lambda^\bullet Q_S\) the parity of exterior degree.

\begin{lemma}[support-Euler selector]
For every positive word \(w\),
\begin{equation}
 \Str\!\left(I\mid\Lambda^\bullet Q_{\supp(w)}\right)
 =(1-1)^{|\supp(w)|-1}=\chi_m(w).
 \label{eq:support-euler}
\end{equation}
The coefficient is invariant under cyclic rotation, relabeling, and
repetition.
\end{lemma}

\begin{proof}
The alternating dimension is

\[
 \sum_{j=0}^{|S|-1}(-1)^j\binom{|S|-1}{j}=(1-1)^{|S|-1}.
\]
For a singleton support, \(\Lambda^0(0)=\C\), so the value is one; otherwise it
is zero.  The asserted invariances follow because only the cardinality of
support occurs.
\end{proof}

Equation \eqref{eq:support-euler} is a genuine positive answer to the
coefficient problem, but it is not yet a stationary transfer representation.
The fiber \(Q_{\supp(w)}\) is selected after the completed word is known.  With
zero differential, the even and odd mixed-support classes cancel in Euler
characteristic; they are not acyclic.  We will not call this virtual
cancellation a homological deletion.

\subsection{A stationary projector character}

Let \(V_m=\C^m\) with basis \(e_1,\ldots,e_m\), and let
\begin{equation}
 P_i=e_ie_i^*,\qquad P_iP_j=\delta_{ij}P_i.
 \label{eq:projectors}
\end{equation}

\begin{proposition}[finite stationary selector]
For \(w=a_{i_1}\cdots a_{i_r}\),
\begin{equation}
 \Tr(P_{i_1}\cdots P_{i_r})=\chi_m(w).
 \label{eq:projector-selector}
\end{equation}
The representation is positive, permutation-natural, and exact at every
repetition.
\end{proposition}

\begin{proof}
If all letters agree, the product is \(P_i\) and has trace one.  If two colors
occur, cyclically rotate the word until an unequal adjacent pair occurs; the
corresponding projector product is zero, and cyclicity of trace gives the
claim.
\end{proof}

The construction is stationary, but its state space already contains one
orthogonal recurrent line per supplied color.  This observation does not by
itself prove minimality: a different nonorthogonal or graded realization
might appear smaller.  The next two sections close precisely that loophole.

\begin{table}[t]
\centering
\small
\caption{The two exact constructions solve different parts of the selector
problem.  Neither yet supplies arithmetic selectivity.}
\label{tab:selector-comparison}
\begin{tabularx}{\textwidth}{L{37mm} c c c Y}
\toprule
Construction & exact words & repetitions & stationary & retained data \\
\midrule
Reduced-support exterior & yes & yes & no & completed support \(S\) \\
Coordinate projectors & yes & yes & yes & one line per supplied color \\
\bottomrule
\end{tabularx}
\end{table}
